% !TeX root = ../../../main.tex
탈출 영역의 가로, 세로 크기 $n$이 주어진다. ($4 \le n \le 1,000$인 정수)
그 다음 줄부터 $n \times n$ 크기의 기룡이가 파악한 탈출 영역이 2차원 배열로 제시된다. 이동 가능한 좌석은 'O', 이동 불가능한 좌석은 'X'로 표현된다. 

탈출이 불가능한 경우는 주어지지 않으며, 탈출 영역의 가장 왼쪽 아래(시작점)와 가장 오른쪽 위(도착점)는 항상 'O'이다.